  \documentclass{article}

  \usepackage{fancyhdr}
  \usepackage{amsmath}
  \usepackage{amsthm}
  \usepackage{amsfonts}
  \usepackage{tikz}
  \usepackage[plain]{algorithm}
  \usepackage{algpseudocode}

  \usetikzlibrary{automata,positioning}

  %
  % Basic Document Settings
  %

  \topmargin=-0.45in
  \evensidemargin=0in
  \oddsidemargin=0in
  \textwidth=6.5in
  \textheight=9.0in
  \headsep=0.25in

  \linespread{1.1}

  \pagestyle{fancy}
  \lhead{\hmwkAuthorName}
  \chead{} % empty
  \rhead{\hmwkHeaderTitle}
  \cfoot{\thepage}

  \renewcommand\headrulewidth{0.4pt}
  \renewcommand\footrulewidth{0pt}

  \setlength\parindent{0pt}

  %
  % Homework Details
  %

  \newcommand{\hmwkTitle}{Homework\ 2}
  \newcommand{\hmwkDueDate}{Tuesday, January 27, 2026,\ 11{:}59 pm (Thai time)}
  \newcommand{\hmwkClass}{Data Structure\ and Abstractions}
  \newcommand{\hmwkClassTime}{Section 2}
  \newcommand{\hmwkClassInstructor}{Kanat Tangwongsan}
  \newcommand{\hmwkAuthorName}{\textbf{Jiraroj Wiruchpongsanon (6781617)}}

  % What appears in the top-right header
  \newcommand{\hmwkHeaderTitle}{Data Structure and Abstractions: Homework 2}

  %
  % Title Page
  %

  \title{
      \vspace{2in}
      \textmd{\textbf{\hmwkClass:\ \hmwkTitle}}\\
      \normalsize\vspace{0.1in}\small{Due\ on\ \hmwkDueDate}\\
      \vspace{0.1in}\large{\textit{\hmwkClassInstructor\ \hmwkClassTime}}
      \vspace{3in}
  }

  \author{\hmwkAuthorName}
  \date{}

  \renewcommand{\part}[1]{\textbf{\large Part \Alph{partCounter}}\stepcounter{partCounter}\\}

  %
  % Helper Commands
  %

  \newcounter{partCounter}
  \newcommand{\solution}{\textbf{\large Solution}}

  \newtheorem*{theorem}{Theorem}

  \theoremstyle{remark}
  \newtheorem*{remark}{Remark}

  \begin{document}

  \maketitle
  \newpage

  \section*{Important Notes from the Assignment}

  \textit{Typesetting Requirements.} For written questions, you must properly typeset your answer,
  produce a PDF file called \texttt{hw2.pdf}, and add it to the zip file that you will upload to Canvas. No
  other format will be accepted. To typeset your homework, apart from Microsoft Word, there are
  LibreOffice and \LaTeX, which we recommend. Note that a scan of your handwritten solution will not be
  accepted. You can find helpful resources about how to typeset your work on the course website.

  \textit{AI Usage.} Use of generative AI (genAI), including AI-assisted code completion, is prohibited.
  You must write every element of the code yourself.

  \textit{Hand-In.} When you hand in your work, your zip file will contain no more than the following
  files: \texttt{hw2.pdf}, \texttt{MinMax.java}, \texttt{RPal.java}, \texttt{mydeque/LinkedListDeque.java},
  and \texttt{mydeque/ArrayDeque.java}.

  \textit{Collaboration.} We interpret collaboration very liberally. You may work with other students.
  However, each student must write up and hand in his or her assignment separately. Let us repeat: You
  need to write your own code. You must not look at or copy someone else's code. You need to write up
  answers to written problems individually. The fact that you can recreate the solution from memory will
  be taken as proof that you actually understood it, and you may actually be interviewed about your
  answers. Be sure to indicate who you have worked with (refer to the hand-in instructions).

  \newpage

  \section*{Problem 2 (8 points)}

  \textbf{Task 2 --- Fibonacci Growth (8 points)}

  For this task, save your work in \texttt{hw2.pdf}.

  In class, our ArrayList implementation uses the array doubling trick, which doubles the capacity of
  the array every time the array becomes full. Our analysis shows that starting with an array of capacity 1
  with no data items initially, appending $n$ data items ends up needing at most $2n$ copying steps.

  This problem involves a more fancy array growing scheme. To describe this new approach, recall
  the Fibonacci sequence from Discrete Math, which is given by
  \[
      F_{n+2} = F_{n+1} + F_n \quad \text{for } n \ge 0,
  \]
  with $F_1 = F_2 = 1$. In our new scheme, the capacity of the underlying array will strictly follow the
  Fibonacci sequence. Initially, the capacity is $F_2 = 1$. When full, we grow the capacity to $F_3 = 2$.
  When full again, we grow the capacity to $F_4 = 3$\ldots{} then to $F_5 = 5$, then to $F_6 = 8$, and so on.

  Amazingly this turns out to work quite well! You will prove a few facts about Fibonacci numbers and
  apply them to analyze the total copying steps.

  \textbf{Subtask I} Using mathematical induction, prove that for $n \ge 1$,
  \[
      1 + F_1 + F_2 + \dots + F_n = F_{n+2}.
  \]
  You are expected to write a solid, rigorous proof based on the definition of Fibonacci numbers (above).

  \textbf{Subtask II} Prove, e.g., using direct proof that for $n \ge 1$,
  \[
      \frac{1}{F_n} \left(1 + \sum_{k=1}^{n} F_k \right) \le 3.
  \]

  \textbf{Subtask III} Suppose we start with no data items in our ArrayList. If we use the Fibonacci growth
  scheme and add in $n$ data items, give a detailed analysis of the total number of copy steps (like we did in
  class) of this scheme. State your argument carefully. To simplify matters, you may wish to assume that
  $n = F_r + 1$ for some integer $r \ge 2$.

  \medskip
  \solution

  \begin{enumerate}
      \item[(I)] 
        \begin{proof}

        For $n \geq 1$ where $n \in \mathbb{I}$, let $S(n) \text{ be } 1 + F_1 + F_2 + ... + F_n = F_{n+2}$; where $F_i$ notates the $i^{th}$ member of a fibonacci sequence, $F_1 = F_2 = 1$ and $F_{n+2} = F_{n+1} + F_n$. \\

        \paragraph{test base case (n = 1,2,3,4).}
          \begin{align*}
          S(1) &= 1 + F_1 = F_{1+2} = 1 + 1 = 2 = F_3,\\
          S(2) &= 1 + F_1 + F_2 = F_{2+2} = 1 + 1 + 1 = 3 = F_4,\\
          S(3) &= 1 + F_1 + F_2 + F_3 = F_{3+2} = 1 + 1 + 1 + 2 = 5 = F_5,\\
          S(4) &= 1 + F_1 + F_2 + F_3 + F_4 = F_{4+2} = 1 + 1 + 1 + 2 + 3 = 8 = F_6.
          \end{align*}  
        So, $S(1), S(2), S(3), S(4)$ holds true.

        % couldn't end subtask I in the page, so enter new page from this part.
        \newpage
        \paragraph{induction}
          we're given $S(n): 1 + F_1 + F_2 + ... + F_n = F_{n+2} \text{ ; } n \geq 1, \text{ and } F_1 = F_2 = 1$. 
        \begin{align*}
        S(n+1)  &= 1 + F_1 + F_2 + ... + F_n + F_{n+1} \\
                &= (1 + F_1 + F_2 + ... + F_n) + F_{n+1}
                && \text{group the first $n^{\text{th}}$ term together} \\
                &= F_{n+2} + F_{n+1}
                && \text{the first term is equivalent to } S(n) \\
                &= F_{n+3}  
                && \text{as the recursive relation of a fibonacci sequence}
        \end{align*}

        \paragraph{conclusion} $S(n)$ holds true for all $n >= 1$

        \qedhere
        \end{proof}

      \item[(II)]
        \begin{proof}
        \begin{align*}
            \text{LHS}
            &= \frac{1}{F_n}\!\left(1 + \sum_{k=1}^n F_k\right) \\
            &= \frac{F_{n+2}}{F_n} 
            && \text{latter term is equivalent to } F_{n+2} \text{ --- proof in Subtask I} \\
            &= \frac{F_{n+1} + F_{n}}{F_n} 
            && \text{unfold $F_{n+2}$ by fibonacci sequence's recursive relation} \\
            &= \frac{F_n + F_{n-1} + F_n}{F_n}
            && \text{unfold $F_{n+1}$ by fibonacci sequence's recursive relation} \\
            &= \frac{2F_n + F_{n-1}}{F_n}
            && \text{use commutative and associative property of $\mathbb{R}$} \\
            &= 2 + \frac{F_{n-1}}{F_n} \\
            &\le 2 + 1 
            &&  F_{n-1} \le F_n \rightarrow \frac{F_{n-1}}{F_n} \le 1\\
            &\le 3 \text{ = RHS.}
        \end{align*}
        \qedhere
        \end{proof}

    % not enough space again
    \newpage

    \item[(III)] We assume that $n = F_r + 1$ where $r \geq 2$ --- as this gives the worse ratio of $n$ to 'number of copies'. Started with $F_2$, when the amount of items exceed our ArrayList size, we copy the the old capacity worth of items, so the total number of copy steps are

        \[F_2 + F_3 + F_4 + ... + F_{r}\]

        which we can use the relation we proof in Subtask I to simplify as $F_{r+2} - 2$ --- worst case steps needed. Now, the result isn't so easy to compare to the 'doubling ArrayList's size approach', so we have to simplify it even more. 

        \begin{align*}
            F_{r+2} - 2 &= F_{r+1} + F_r - 2
                && \text{unfold $F_{r+2}$ by fibonacci sequence's recursive relation} \\
                &= F_r + F_{r-1} + F_r - 2
                && \text{unfold $F_{r+1}$ by fibonacci sequence's recursive relation} \\
                &= 2n + F_{r-1} - 4
                && \text{since } n = F_r + 1 \\
                & \le 3n - 5
                &&  \text{since } F_r = n - 1, \text{ and } F_{n-1} \le F_n \rightarrow \frac{F_{n-1}}{F_n} \le 1\\
                & \le 3n
        \end{align*}

        So, it turns out using a fibonacci sequence as ArrayList size incrementor performs worse than just doubling the ArrayList's size; this probably due to the sequence's growth scheme being less than 2 --- as defined the the $n^{th}$ fibonacci number as $F_n = F_{n-1} + F_{n-2}$ where $n \geq 1$; $F_{n-2} \le F_{n-1} \rightarrow F_{n} \le 2F_{n-1}$

  \end{enumerate}

  \end{document}
